\chapter{Dataset and Experimental Protocol}
\label{ch:protocol}

Having derived and implemented eALS, we now turn to how it is evaluated:
which dataset was used and how it was filtered and split, which metrics
quantify accuracy, and which hyperparameters are held fixed across the
experiments of Chapters~\ref{ch:results} and~\ref{ch:scalability}.

\section{Dataset}

The paper evaluates on two datasets, the Yelp Challenge dataset and Amazon
``Movies and TV'' reviews, and the assignment brief asks for the
experimental analysis to use one of them rather than a substitute such as
MovieLens, so the choice between the two was really a choice of which one
was still reachable. The Yelp Challenge dataset, in the specific form the
paper used it, a static 2015 snapshot of roughly 1.6 million reviews
distributed directly by Yelp as a downloadable archive for that
competition, is no longer available. Yelp has since folded that data into
an ``Open Dataset'' aimed at a different use case (business listings for
their own developer program): not the same snapshot, not the same schema,
and not redistributed in a form that would let this project reproduce the
paper's exact filtering pipeline. Reconstructing something merely
\emph{similar} to it would have quietly broken the assignment's requirement
of testing on ``one of the datasets used by the paper'', so this project
uses the Amazon dataset instead, which remains available essentially as
the paper used it: the ratings-only CSV hosted by the Stanford Network
Analysis Project \citep{amazonreviews}, downloaded directly and not
modified beyond the filtering described below. It contains
4{,}607{,}047 raw \texttt{(user, item, rating, timestamp)} rows. Every
rating is binarized to a positive implicit interaction ($r_{ui}=1$); the
rating value itself is discarded, matching the paper's treatment of the
data (``we transform the review dataset into implicit data, where each
entry is marked as 0/1 indicating whether the user reviewed the item'').

\subsection{Filtering}

Following the paper, users and items with fewer than 10 interactions are
removed, since extremely sparse users and items make offline evaluation
unreliable (over half the users in the raw Amazon data have only a single
review) without being representative of what a production recommender
would face. Removing sparse users can push some items below the threshold
and vice versa, so the filter has to be applied \emph{iteratively} until it
reaches a fixed point. \texttt{src/preprocess.py} repeats the
filter-and-recount step up to 50 times and explicitly checks for
convergence rather than assuming a fixed number of passes is enough. On
this dataset the filter took 20 rounds to fully stabilize (from
4{,}607{,}047 interactions after the first pass down to 1{,}204{,}688, then
shrinking by a decreasing amount each round until round 20 leaves the count
unchanged at 958{,}986). An earlier version of this pipeline capped the
filter at 10 rounds, which removes over 99.9\% of what the full filter
would remove but is not exact; this was caught and corrected, and
Table~\ref{tab:dataset} reports the fully converged numbers.

\begin{table}[H]
\centering
\begin{tabular}{lrr}
\toprule
 & Paper (Amazon) & This project \\
\midrule
Users        & 117{,}176   & 33{,}326 \\
Items        & 75{,}389    & 21{,}901 \\
Interactions & 5{,}020{,}705 & 958{,}986 \\
Sparsity     & 99.94\%     & 99.87\% \\
\bottomrule
\end{tabular}
\caption{Dataset statistics after 10-core filtering, paper versus this
project.}
\label{tab:dataset}
\end{table}

\medskip
The two datasets are visibly not the same scale. This is not a filtering
bug: the CSV currently hosted by SNAP is a later, differently-curated
snapshot of Amazon review data than the one the authors filtered in 2016,
so the raw interaction counts already diverge before any filtering is
applied. Section~\ref{sec:weaknesses} returns to what this scale
difference does and does not change about the conclusions drawn from the
experiments.

\subsection{Splitting}

Interactions are re-indexed to dense integer ranges $[0,M)$ and $[0,N)$ and
split with the same leave-one-out protocol as the paper's offline setting
(Section 5.1 of \citealp{he2016eals}): each user's chronologically latest
interaction is held out as the test item, and the remaining interactions
form the training set. This protocol evaluates a model's ability to
predict a known user's next interaction from their history; it does not
test cold-start behaviour for new users, since by construction every test
user has training history.

\section{Evaluation metrics}

Two ranking metrics are used, matching the paper: \textbf{Hit Ratio}
(HR@$k$), the fraction of test users for whom the held-out item appears in
the top-$k$ of the model's ranked recommendations, and \textbf{Normalized
Discounted Cumulative Gain} (NDCG@$k$), which additionally rewards ranking
the held-out item near the top rather than merely inside the top-$k$:
\[
  \mathrm{NDCG@}k = \frac{1}{|\mathcal{T}|}\sum_{(u,i)\in\mathcal{T}}
  \frac{\mathbb{1}[\mathrm{rank}(u,i) < k]}{\log_2(\mathrm{rank}(u,i)+2)},
\]
where $\mathcal{T}$ is the test set and $\mathrm{rank}(u,i)$ is the
position of the held-out item $i$ in $u$'s ranked list (0-indexed). Both
metrics are implemented in \texttt{src/evaluate.py}.

\medskip
One deliberate deviation from the paper is worth stating plainly here
rather than leaving it implicit. The paper ranks the held-out item against
the \emph{entire} item catalog for every test user. This is the exact
protocol, and \texttt{evaluate\_full} in \texttt{src/evaluate.py}
implements it faithfully; it is exercised in the small-data notebook
(Appendix~\ref{app:notebook}), where re-ranking the whole 45-item toy
catalog costs almost nothing. For the large-scale experiments in
Chapters~\ref{ch:results} and~\ref{ch:scalability}, which retrain the model
dozens of times over a grid of hyperparameters, ranking against the full
21{,}901-item catalog for all 33{,}326 test users on every single
evaluation point would dominate the total experiment time and add nothing
the scalability study is meant to measure. Instead,
\texttt{evaluate\_sampled} ranks the held-out item against 99 items the
user never interacted with, sampled uniformly at random, and reports
HR@10/NDCG@10 on that shorter list. This is the same relaxation used in
later work by the same first author \citep{he2017ncf}, and it is now
standard practice in the recommender-systems literature for exactly this
reason.

\section{Experimental setup}

Unless stated otherwise, experiments use $\lambda=0.01$ (matching the
paper's choice, made ``for a fair comparison'' across methods),
$w_{ui}=1$ for observed interactions (also matching the paper), and
$K=32$ latent factors as the default operating point, with dedicated
sweeps over $K \in \{8,16,32,64,128\}$ for the scalability study. All
experiments run on a single machine (Spark in local mode, up to 8
partitions unless the experiment is specifically varying the partition
count) rather than a real cluster; Chapter~\ref{ch:discussion} discusses
what this does and does not tell us about behaviour on an actual
distributed deployment.
